%% kapa: exam=1 round=30 year=2019 period=1 form=A subject=economics subject_ko=경제학원론 seq=27 number=67
%% source: Q-Net 감정평가사 기출문제 (cst00309) artlSeq=5205205 file=extracted/1008/5205205/2019년 제30회 감정평가사 1차시험 문제지/2019년 제30회 감정평가사 1차시험 문제지/2019년도 제30회 감정평가사 1차 1교시 A형.hwp
%% conversion: kordoc 4.24.51 -> tools/gen1.py (auto); review=proofread-sonnet; flags=none
\begin{question}{67}
甲국의 총생산함수가 ${ Y=AK^{ 0.4 } L^{ 0.6 } }$이다. 甲국 경제에 관한 설명으로 옳은 것을 모두 고른 것은? (단, ${ Y }$는 생산량, ${ A }$는 총요소생산성, ${ K }$는 자본량, ${ L }$은 노동량으로 인구와 같다.)
\begin{framedtext}
ㄱ. 생산량의 변화율을 노동량의 변화율로 나눈 값은 0.6으로 일정하다. \par
ㄴ. ${ A }$가 3 \% 증가하면, 노동의 한계생산도 3 \% 증가한다. \par
ㄷ. 1인당 자본량이 2 \% 증가하면, 노동의 한계생산은 1.2 \% 증가한다. \par
ㄹ. ${ A }$는 2 \% 증가하고 인구가 2 \% 감소하면, 1인당 생산량은 2.8 \% 증가한다.
\end{framedtext}
\begin{choices}
  \item ㄱ, ㄹ
  \item ㄴ, ㄷ
  \item ㄷ, ㄹ
  \item ㄱ, ㄴ, ㄹ
  \item ㄱ, ㄷ, ㄹ
\end{choices}
\end{question}
\answer{4}
